\documentclass[12pt,a4paper]{article}

% --- Packages ---
\usepackage[utf8]{inputenc}
\usepackage[T5]{fontenc}
\usepackage[vietnamese,english]{babel}
\usepackage{geometry}
\geometry{margin=2.5cm, top=3cm, bottom=3cm}
\usepackage{amsmath,amssymb}
\usepackage{graphicx}
\usepackage{booktabs}
\usepackage{array}
\usepackage{multirow}
\usepackage{xcolor}
\usepackage{hyperref}
\usepackage{caption}
\usepackage{subcaption}
\usepackage{float}
\usepackage{enumitem}
\usepackage{setspace}
\usepackage{listings}

\lstset{
  basicstyle=\ttfamily\small,
  breaklines=true,
  showstringspaces=false,
  frame=single,
  framesep=4pt,
  columns=fullflexible,
}

\hypersetup{
  colorlinks=true,
  linkcolor=blue!60!black,
  urlcolor=blue!60!black,
  citecolor=red!60!black,
}

\title{\textbf{So sánh hai implementation tái hiện bài báo SAC:\\
\emph{Semantic-Aware-Video-Compression-for-Automotive-Cameras}\\
\textbf{vs.} \emph{new\_branch}}}
\author{Nguyễn Hữu Huy}
\date{17 tháng 5, 2026}

\begin{document}
\maketitle

\section{Mục tiêu báo cáo}

Cả hai folder cùng nhằm tái hiện bài báo \emph{Semantic-Aware Video Compression for Automotive
Cameras} (Wang et al., IEEE TIV 2023):
\begin{itemize}[leftmargin=2em]
  \item \texttt{scripts/Semantic-Aware-Video-Compression-for-Automotive-Cameras/}
        — viết tắt \textbf{[paper-folder]}. Đây là fork đã được vá lỗi từ repo gốc của
        nhóm tác giả, dùng kiến trúc \emph{tự thiết kế} (không phải CCNet chuẩn).
  \item \texttt{scripts/new\_branch/} — viết tắt \textbf{[new-branch]}. Đây là bản tái
        hiện ``đúng nguyên văn CLAUDE.md'' — dùng CCNet ResNet-101 dilated + RCCA gốc.
\end{itemize}

Mặc dù \emph{new-branch} sát với mô tả paper hơn về mặt câu chữ, \textbf{kết quả thực
tế của paper-folder lại tốt hơn} (tiệm cận con số trong Bảng II/IV gốc). Báo cáo này
trả lời ba câu hỏi:
\begin{enumerate}[leftmargin=2em]
  \item Sự khác biệt cụ thể giữa hai bản code là gì?
  \item Nguyên nhân nào gây ra chênh lệch kết quả?
  \item Có thể cải thiện kết quả của bản \emph{code-paper} (new-branch) như thế nào?
\end{enumerate}

\section{Bảng tổng quan các khác biệt}

\begin{table}[H]
\centering
\small
\caption{So sánh side-by-side toàn pipeline. Ô \textbf{in đậm} là phần khác biệt cốt
lõi đóng vai trò quan trọng nhất cho chênh lệch kết quả.}
\label{tab:overall}
\begin{tabular}{p{3.3cm} p{5.4cm} p{5.4cm}}
\toprule
\textbf{Hạng mục} & \textbf{paper-folder} & \textbf{new-branch} \\
\midrule
Kiến trúc backbone &
  CNN tự thiết kế: 3 res-block 64$\to$128$\to$256 + 3 down (MaxPool), output $512$ ch &
  \textbf{ResNet-101 dilated} chuẩn CCNet (4 layer, output $2048$ ch, stride 8) \\
\addlinespace
Normalization &
  \textbf{GroupNorm(4, $\cdot$)} + LeakyReLU &
  BatchNorm2d + ReLU (CCNet gốc thay từ InPlaceABNSync) \\
\addlinespace
Tham số &
  $\sim$10--15 M &
  $\sim$70 M (ResNet-101 + RCCA + DSN) \\
\addlinespace
Decoder head &
  \textbf{Multi-scale fusion}: RCCA $\to$ bottleneck + 3 nhánh ConvTranspose (S3,S2,S1)
  concat $\to$ out conv &
  Bottleneck đơn $\to$ bilinear $\times 8$ \\
\addlinespace
Aux supervision &
  \emph{Không có} (đã có multi-scale ở head) &
  \textbf{DSN} từ layer3, loss $= L_{\text{main}} + 0.4\,L_{\text{aux}}$ \\
\addlinespace
Output &
  \texttt{softmax(logits)} trả về trong forward &
  Logits thô; softmax bên ngoài khi tính dice \\
\addlinespace
Loss &
  Dice + soft cross-entropy với one-hot label đã làm mịn &
  Dice + \texttt{nn.CrossEntropyLoss} trên label rắn 0--3 \\
\addlinespace
Pretrain &
  \emph{Không}, train from scratch &
  ResNet-101 ImageNet pretrained \\
\addlinespace
Tiền xử lý ảnh &
  \textbf{RGB} (3 kênh gốc), resize bicubic order=3, \textbf{z-score per-image} &
  \textbf{Grayscale}-as-3ch (Y luma), resize bilinear, chuẩn hoá ImageNet mean/std \\
\addlinespace
Label resize &
  \textbf{One-hot $\to$ resize order=1 (linear)} = label mềm ở biên &
  Nearest-neighbor trên class id rắn \\
\addlinespace
Augment &
  \emph{Không} &
  Random horizontal flip + random scale (0.75--1.25) + random crop \\
\addlinespace
Optimizer & Adam 3e-4, $\beta=(0.9,0.999)$, wd 1e-5, batch 4, FP16 &
                Adam 3e-4 idem, batch 4, AMP, clip-grad 35 \\
\addlinespace
Epochs & 200 (mặc định), checkpoint mỗi 50 step & 47 (theo paper) \\
\midrule
\textbf{Khối S2 — sinh mask ROI} &
  \textbf{Dùng GT $\,$\texttt{\_gtFine\_4class.png}} cho cả train/val/test (đã merge trực tiếp từ Cityscapes GT) &
  \textbf{Dùng prediction} của mạng segmentation trên ảnh gốc \\
\addlinespace
Macroblock filter &
  Vòng for Python, kiểm tra \texttt{patch.max()==1} &
  \texttt{F.max\_pool2d(kernel=16, stride=16)} + upsample (vectorised) \\
\addlinespace
Codec params &
  \texttt{libx264}/\texttt{libx265}, FPS=\textbf{17}, preset medium, không cố định GOP &
  Idem, FPS=\textbf{30}, GOP cố định \texttt{keyint=min-keyint=30} \\
\addlinespace
Reconstruct &
  \texttt{cv2.add(IA, BA)} — saturating add &
  \texttt{np.where(mask, roi\_dec, non\_dec)} dùng lại mask gốc \\
\addlinespace
iIOU đánh giá &
  Binary GT vs binary pred (file PNG đã lưu) &
  Tính mIoU 4 lớp + iIoU lớp ROI từ pred (sau khi seg lại trên ảnh tái tạo) \\
\bottomrule
\end{tabular}
\end{table}

\section{Phân tích nguyên nhân paper-folder cho kết quả tốt hơn}

\subsection{Nguyên nhân \#1 — Khối S2 dùng MASK GT thay vì PRED (quan trọng nhất)}

\textbf{Đây là chênh lệch lớn nhất.} Đọc trực tiếp
\texttt{Codes/CCNet/TwoStream\_generate.py} dòng 84:
\begin{lstlisting}[language=Python]
seg_path = _get_label_path(img_path)        # _gtFine_4class.png
seg_4class = cv2.imread(seg_path, cv2.IMREAD_GRAYSCALE)
...
roi_binary = (seg_4class == 0).astype(np.uint8)
roi_macro  = seg_mask_im_macro(roi_binary)
\end{lstlisting}

paper-folder \emph{không dùng output của mạng} để tạo mask. Mask ROI lấy từ
\textbf{ground-truth 4 lớp} đã merge sẵn của Cityscapes. Hệ quả:
\begin{itemize}[leftmargin=2em]
  \item Không có lỗi segmentation lan sang khối nén $\to$ ROI thật sự được phủ chính
        xác, không bỏ sót pixel vật thể quan trọng.
  \item \texttt{iIOU} sau nén được đo bằng cách so GT 4-class với mask
        \emph{dự đoán bởi mạng trên frame tái tạo} (\texttt{outputs/ccnet\_test\_debug}).
        Vì ROI input đã được nén nhẹ (CRF=18) bằng mask GT, model dễ dàng đạt iIOU cao.
\end{itemize}

Trong khi đó, new-branch (\texttt{evaluate.py} dòng 174):
\begin{lstlisting}[language=Python]
seg_lowres = _seg_predict(model, img_t, device)
roi_masks.append(build_roi_mask(seg_lowres, target_hw=FRAME_HW))
\end{lstlisting}

dùng prediction. Nếu mạng segmentation chưa đạt iIoU $>$ 95\%, một số pixel ROI sẽ bị
gán nhãn non-ROI, bị nén CRF=27, mất chi tiết, kéo mIoU \emph{sau} nén xuống.

\textbf{Hiệu ứng định lượng (ước tính):} chỉ cần $\sim$5\% pixel ROI bị mis-classify
thành non-ROI ở biên, mIoU sau nén có thể giảm $1.5$--$3$ điểm — đủ để new-branch ra
con số $\approx$ 87\% thay vì 90.56\% như paper.

\subsection{Nguyên nhân \#2 — Multi-scale decoder vs upsample đơn}

paper-folder dùng đầu ra phân vùng \emph{tổng hợp} 3 mức (xem
\texttt{Codes/CCNet/CCNet.py} dòng 136--163):
\begin{lstlisting}[language=Python]
self.S3 = nn.ConvTranspose2d(512, 4, kernel_size=(8,8), stride=(8,8))  # 1/8 -> 1/1
self.S2 = nn.ConvTranspose2d(128, 4, kernel_size=(4,4), stride=(4,4))  # 1/4 -> 1/1
self.S1 = nn.Conv2d(32, 4, kernel_size=1)                              # 1/1
S = torch.cat([S3, S2, S1], dim=1)        # 12 channels
output = self.out(S)
\end{lstlisting}

Ba dự đoán ở ba độ phân giải được nối lại rồi fuse. Đây là \emph{multi-scale
prediction fusion} — kỹ thuật chuẩn giúp:
\begin{itemize}[leftmargin=2em]
  \item Lớp khối lớn (\texttt{sky}, \texttt{construction}) học tốt từ S3 (semantic
        đậm, độ phân giải thấp).
  \item Biên/chi tiết nhỏ (cột đèn, biển báo, người đi bộ) học tốt từ S1 (texture giữ
        nguyên ở mức pixel).
\end{itemize}

new-branch chỉ có một head đầu ra duy nhất rồi bilinear-upsample $\times 8$:
biên không sắc nét, dễ ``rách'' khi macroblock 16$\times$16 mở rộng vùng ROI sai.

\subsection{Nguyên nhân \#3 — GroupNorm thay BatchNorm}

Cityscapes train có $\sim$2975 ảnh, batch=4. Với BN, mỗi mini-batch $4$ ảnh đô-thị
khác nhau (4 thành phố) cho thống kê rất nhiễu. GroupNorm độc lập batch, ổn định hơn
nhiều ở batch nhỏ — đặc biệt quan trọng cho mạng \emph{tự thiết kế nhỏ} có chiều rộng
hạn chế, không nhờ được pretrain trên ImageNet làm ``mỏ neo''.

new-branch dùng BN cho cả ResNet-101 lẫn RCCA module — phải bù bằng ImageNet pretrain
(nên các thống kê BN bắt đầu hợp lý). Nhưng trong giai đoạn đầu fine-tune, BN
running\_mean/var bị shift, gây ra dao động dữ dội ở $\sim$5 epoch đầu.

\subsection{Nguyên nhân \#4 — Label mềm tại biên (label smoothing tự nhiên)}

paper-folder \texttt{read.py} (\texttt{resize\_label}):
\begin{lstlisting}[language=Python]
result[i] = resize(image[i].astype(float),
                   (edsize[1]//2, edsize[2]//2),
                   order=1, mode='edge')          # linear, KHONG nearest
\end{lstlisting}

Label được \emph{one-hot trước khi resize} rồi nội suy tuyến tính. Pixel ở biên giữa
ROI/sky không còn ``0 hoặc 1'' mà là $\sim 0.3$/$0.7$. Khi dùng với
\texttt{soft\_cross\_entropy(y\_true, log(y\_pred))}, đây chính là \textbf{label
smoothing có cấu trúc}. Hiệu ứng:
\begin{itemize}[leftmargin=2em]
  \item Loss không phạt nặng dự đoán ``0.5'' ở biên — model học mềm dẻo hơn.
  \item Giảm overfitting trên biên răng cưa của ground-truth.
\end{itemize}

new-branch dùng \texttt{Image.NEAREST} resize label rắn $\to$ \texttt{CrossEntropyLoss}
phạt rất mạnh ở biên (mỗi pixel chỉ có 1 nhãn đúng). Đặc biệt ở pixel ``cây gần xe'',
``cột đèn dọc theo tường'' — biên ground-truth bản thân đã không hoàn hảo.

\subsection{Nguyên nhân \#5 — Per-image z-score vs ImageNet stats}

paper-folder chuẩn hoá theo chính mỗi ảnh:
\begin{equation}
\hat{x} = \frac{x - \mu_x}{\sigma_x + 10^{-8}}
\end{equation}

Phương pháp này \emph{robust với illumination} (ngày nắng, hoàng hôn, đường hầm). Vì
model train from-scratch nên không cần khớp pretrained-stats. Cityscapes vốn chụp ở
nhiều điều kiện ánh sáng, nên per-image z-score bù trừ tốt.

new-branch chuẩn hoá theo ImageNet $(\mu, \sigma)$ = $(0.485, 0.456, 0.406)$ /
$(0.229,0.224,0.225)$ — đúng cho pretrained ResNet-101 nhưng \emph{không bù trừ}
biến thiên sáng tối. Thêm vào đó, new-branch \emph{convert sang grayscale} (theo
CLAUDE.md \S 5.1) làm mất hoàn toàn thông tin màu — đặc biệt với 4 lớp mà \texttt{sky}
phụ thuộc mạnh vào kênh blue, \texttt{nature} phụ thuộc kênh green, đây là điểm yếu
rõ ràng.

\subsection{Nguyên nhân \#6 — Codec config (FPS=17 đúng với Cityscapes)}

Cityscapes \texttt{leftImg8bit\_sequence} thực ra ghi ở $\sim$17 fps (camera xe).
paper-folder dùng \texttt{-framerate 17}; new-branch dùng \texttt{-framerate 30}. Ở
30 fps, codec nghĩ chuyển động chậm hơn thật $\to$ inter-frame prediction kém hiệu
quả $\to$ cùng CRF nhưng \emph{bitrate} cao hơn $\to$ chất lượng cao hơn $\to$
\emph{lợi cho new-branch về PSNR/SSIM}.

Tuy vậy, khi đo \emph{cùng tỉ lệ nén} (paper Bảng IV ép cùng compression ratio),
new-branch sẽ kém hơn vì 30 fps che lấp lợi thế semantic-aware. \textbf{Để công bằng
với paper, new-branch nên dùng FPS=17}.

Ngoài ra, paper-folder không ép \texttt{keyint}, để x264 chọn GOP tự động (mặc định
$\sim$250 frame). new-branch ép \texttt{keyint=30} cho mỗi GOP, khiến mỗi giây có 1
I-frame — quá nhiều, tăng bitrate không cần thiết.

\subsection{Nguyên nhân \#7 — Auxiliary head có ích hay không?}

new-branch có DSN aux loss ($\times 0.4$). Trên paper gốc CCNet (Cityscapes 19 lớp),
DSN \emph{có giúp} $\sim$1 điểm mIoU. Tuy nhiên với chỉ \textbf{4 lớp gộp}, layer3
features đã đủ rich, DSN không thêm gì nhiều mà còn:
\begin{itemize}[leftmargin=2em]
  \item Tăng chi phí FLOPs khi training.
  \item Gradient từ DSN có thể ``kéo'' feature theo hướng tối ưu cho dự đoán thô
        $\to$ làm khối main head hơi chậm hội tụ.
\end{itemize}

paper-folder thay thế bằng \emph{multi-scale prediction fusion} (S3/S2/S1) — gradient
chảy về cả 3 đầu nhưng đều cộng dồn vào output cuối $\to$ supervised dense, hợp lý
hơn cho task 4 lớp.

\subsection{Tóm tắt mức độ đóng góp}

\begin{table}[H]
\centering
\caption{Ước lượng đóng góp tương đối của từng nguyên nhân tới khoảng cách kết quả.}
\label{tab:contrib}
\begin{tabular}{lcc}
\toprule
Nguyên nhân & Hướng & Ước lượng tác động \\
\midrule
\#1 Dùng GT mask thay vì pred & paper $\uparrow$ & \textbf{rất lớn} ($\sim$2--3 điểm mIoU sau nén) \\
\#2 Multi-scale decoder       & paper $\uparrow$ & lớn ($\sim$0.5--1 điểm mIoU base) \\
\#3 GroupNorm vs BN, batch 4  & paper $\uparrow$ & vừa ($\sim$0.3--0.5 điểm) \\
\#4 Label mềm khi resize      & paper $\uparrow$ & vừa ($\sim$0.3 điểm + giảm artefact biên) \\
\#5 Per-image z-score + giữ RGB (combo) & paper $\uparrow$ & lớn ($\sim$0.5--1 điểm; chỉ có hiệu lực nếu \textbf{re-train}) \\
\#6 FPS 17 + GOP dài (combo)  & paper $\uparrow$ & nhỏ-vừa với \texttt{leftImg8bit\_sequence}; \textbf{phản tác dụng} với slideshow eval (xem \S\ref{sec:pitfall}) \\
\#7 Có DSN aux                & new-branch $\uparrow$ (nhẹ) & nhỏ \\
\bottomrule
\end{tabular}
\end{table}

\section{Đề xuất cải thiện code paper (new-branch)}

Mục tiêu: giữ kiến trúc CCNet ResNet-101 đúng paper, nhưng đạt mIoU sau nén
$\ge 90\%$. Đề xuất theo độ ưu tiên giảm dần:

\subsection{Đề xuất 1 — Tách 2 chế độ đánh giá: oracle (GT mask) vs end-to-end (pred mask)}

\textbf{Lý do.} Bảng II của paper cốt yếu đo \emph{compression quality} (PSNR, SSIM,
SA-*) — không nhất thiết cần mạng segmentation. Bảng IV mới đo \emph{segmentation
sau nén} (mIoU, iIoU). Để so sánh trực tiếp với paper, new-branch cần thêm chế độ
\textbf{oracle}: lấy mask ROI từ GT \texttt{\_gtFine\_4class.png}, không qua mạng.

\textbf{Cài đặt.} Trong \texttt{evaluate.py}, thêm flag:
\begin{lstlisting}[language=Python]
ap.add_argument("--roi-source", choices=["pred", "gt"], default="pred")
...
if args.roi_source == "gt":
    gt_4class = np.array(Image.open(lbl_path))
    if gt_4class.shape[:2] != FRAME_HW:
        gt_4class = np.array(Image.fromarray(gt_4class)
                             .resize((FRAME_HW[1], FRAME_HW[0]),
                                     Image.NEAREST))
    seg_lowres = gt_4class[::2, ::2]   # downsize 1:2 cho seg_input_hw
else:
    seg_lowres = _seg_predict(model, img_t, device)
\end{lstlisting}

Sau đó chạy:
\begin{lstlisting}[language=bash]
python -m scripts.new_branch.evaluate --split test --roi-source gt
\end{lstlisting}

Khi đó kết quả sẽ \emph{ngang ngửa} paper-folder. Đây là phép so sánh \textbf{fair vs
paper} vì paper-folder ngầm dùng GT.

\subsection{Đề xuất 2 — Khắc phục label resize cứng (nearest)}

\textbf{Lý do.} Giống Nguyên nhân \#4 ở trên. Có 2 cách:
\begin{enumerate}[leftmargin=2em]
  \item \textbf{One-hot resize bilinear} (giống paper-folder): biến target thành
        $(C,H,W)$ float, resize bilinear, truyền vào dice (mềm) và soft-CE
        ($-\sum y\log\hat{p}$).
  \item \textbf{Edge-aware label smoothing}: giữ class id rắn, nhưng đánh nhãn
        \texttt{ignore\_index=255} cho pixel sát biên (dilate biên 2--3 px). CCNet
        gốc thường làm cách này.
\end{enumerate}

Khuyến nghị cách (1) — gần với paper hơn:
\begin{lstlisting}[language=Python]
# trong __getitem__:
lbl_full = lbl_full.resize((W*2, H*2), Image.NEAREST)   # tinh
lbl_oh = F.one_hot(torch.from_numpy(np.array(lbl_full)).long(), 4) \
            .permute(2,0,1).float()
lbl_soft = F.interpolate(lbl_oh.unsqueeze(0), size=(H,W),
                         mode='bilinear', align_corners=False).squeeze(0)
\end{lstlisting}

và đổi loss thành:
\begin{lstlisting}[language=Python]
ce = -(lbl_soft * F.log_softmax(logits, dim=1)).sum(1).mean()
\end{lstlisting}

\subsection{Đề xuất 3 — Re-train với RGB + augment mạnh + cosine LR}

\textbf{Tiền đề.} Phải sửa \texttt{dataset.py} bỏ \texttt{\_rgb\_to\_gray3} thì
checkpoint mới mới hợp lệ với input RGB. \emph{Không} bật riêng cờ
\texttt{--no-gray} trong \texttt{evaluate.py} với checkpoint cũ — sẽ tụt mIoU
xuống $\sim$50\% (xem \S\ref{sec:pitfall} để biết chi tiết).

\textbf{Lý do giữ RGB.} CLAUDE.md \S 5.1 trích ``\ldots transformed to greyscale''
— nhưng đọc kỹ paper Section III-B.1 thì câu này mô tả \emph{pipeline tổng
quát}, không bắt buộc input mạng phải grayscale. Mất kênh màu khiến \texttt{sky}
(blue) và \texttt{nature} (green) khó phân biệt khi luma giống nhau (ví dụ trời
u ám = lá cây khô).

\textbf{Combo huấn luyện đề xuất:}
\begin{itemize}[leftmargin=2em]
  \item \textbf{(a) Bỏ grayscale} trong \texttt{dataset.py}: thay
        \texttt{img\_tensor = \_normalize(\_rgb\_to\_gray3(img\_np))} bằng
        \texttt{img\_tensor = \_normalize(img\_np)}.
  \item \textbf{(b) Augment mạnh hơn:} random scale 0.5--2.0 (rộng hơn 0.75--1.25
        hiện tại), color jitter (brightness/contrast/saturation $\pm 0.3$),
        random crop $512\times 1024$.
  \item \textbf{(c) Cosine LR schedule:} \texttt{CosineAnnealingLR(T\_max=epochs)}
        thay constant 3e-4. Cosine từ 3e-4 $\to$ 1e-6 trong 80 epoch thường cho
        $+0.5$ điểm mIoU trên Cityscapes.
  \item \textbf{(d) Warmup} 1 epoch đầu (poly hoặc linear) để BN ổn định.
  \item \textbf{(e) Train 80 epoch} (gấp đôi paper) — RTX 4090 mất $\sim$8 giờ.
\end{itemize}

Khi đã có checkpoint mới (\emph{RGB-trained}), eval bằng:
\begin{lstlisting}[language=bash]
python -m scripts.new_branch.evaluate --split test --roi-source gt --no-gray
\end{lstlisting}

\subsection{Đề xuất 4 — Đổi BN trong RCCA sang GN khi train batch nhỏ}

Trong \texttt{scripts/CCNet/networks/ccnet.py} chỉ thay class \texttt{InPlaceABNSync}:
\begin{lstlisting}[language=Python]
class InPlaceABNSync(nn.Module):
    def __init__(self, num_features, **kwargs):
        super().__init__()
        # 32 group la chuan, neu num_features < 32 thi lay min(32, n_feat//4)
        g = max(1, min(32, num_features // 4))
        self.bn = nn.GroupNorm(g, num_features)
        self.act = nn.ReLU(inplace=True)
    def forward(self, x): return self.act(self.bn(x))
\end{lstlisting}

\emph{Chú ý:} ResNet-101 pretrained có sẵn BN running\_stats; chuyển sang GN sẽ
\emph{huỷ} pretrained-stats cho phần BN. Cần re-finetune kỹ hơn ($\ge$ 60 epoch). Vì
vậy đây là đề xuất \textbf{phụ}, chỉ áp dụng khi batch $\le 4$.

\subsection{Đề xuất 5 — Multi-scale prediction fusion thêm vào RCCA head}

Bổ sung 1 head phụ ở stride 4 (lấy feature từ layer2):
\begin{lstlisting}[language=Python]
# Trong networks/ccnet.py ResNet.__init__:
self.aux_s4 = nn.Sequential(
    nn.Conv2d(512, 128, 3, padding=1), InPlaceABNSync(128),
    nn.Conv2d(128, num_classes, 1))

# forward:
x = self.layer1(x); x_aux_s4 = self.aux_s4(x_l2_features)
...
return [x_main, x_dsn, x_aux_s4]   # 3 head, weights = 1.0, 0.4, 0.3
\end{lstlisting}

Ý tưởng bắt chước S1/S2/S3 của paper-folder nhưng giữ kiến trúc ResNet-101.

\subsection{Đề xuất 6 — Đánh giá theo cùng compression-ratio (bitrate khớp)}

Paper Bảng IV ép tất cả phương pháp về \emph{cùng tỉ lệ nén}; nếu chỉ cố định CRF,
SA-X264 thường ra dung lượng \emph{nhỏ hơn} H.264 baseline $\to$ so sánh không
fair. Nên:
\begin{enumerate}[leftmargin=2em]
  \item Đo bytes của H.264 (CRF=23) làm baseline.
  \item Hiệu chỉnh \texttt{crf\_non} của SA-X264 sao cho tổng bytes (IA+BA) khớp
        baseline trong sai số $\pm 2\%$. Có thể chạy binary search trên
        \texttt{crf\_non} $\in [25, 35]$.
\end{enumerate}

Kết quả sau khi đã \emph{khớp bitrate} mới là phép so sánh hợp lệ với paper.

\subsection{Đề xuất 7 — Khi tính iIOU sau nén, tách 2 loại: ROI-GT-IoU và ROI-pred-IoU}

Hiện \texttt{metrics.iiou} so prediction sau nén với GT lớp ROI. Nên thêm cột:
\begin{itemize}[leftmargin=2em]
  \item \texttt{iIoU\_pred-vs-GT} (đang có) — đánh giá tổng thể.
  \item \texttt{iIoU\_pred\_after-vs-pred\_before} — đo riêng \emph{drop} do nén,
        loại bỏ ảnh hưởng của lỗi seg ban đầu. Đây là chỉ số \emph{thực sự đo tác
        dụng của SAC vs baseline}.
\end{itemize}

\subsection{Đề xuất 8 — Dùng leftImg8bit\_sequence thật + codec config phù hợp}

\textbf{Vấn đề.} Hiện \texttt{evaluate.py} chia test thành GOP 30 frame nhưng
\emph{30 frame đó là 30 ảnh từ 30 thành phố/30 scene khác nhau}, không phải đoạn
video liên tục $\to$ codec phải mã hoá như slideshow, mỗi frame gần như I-frame
$\to$ SAC không tận dụng được inter-prediction. Đây là lý do hiện tại phải đặt
\texttt{--keyint 30 --framerate 30} (đảm bảo mỗi frame = 1 I-frame, không bị
codec ép P-frame qua scene cut). Nếu đặt \texttt{--keyint 0} với slideshow, codec
fail $\to$ bytes/frame \emph{tăng} $\sim$25\% (xem \S\ref{sec:pitfall}).

\textbf{Cách đúng:} Tải \texttt{leftImg8bit\_sequence\_trainvaltest.zip} (324 GB)
— mỗi sequence là 30 frame liên tục cùng scene. Cho mỗi sequence:
\begin{enumerate}[leftmargin=2em]
  \item Predict seg cho frame thứ 20 (frame có nhãn).
  \item Dùng cùng mask đó cho toàn bộ 30 frame (giả định scene ít thay đổi trong
        $\sim$2s, hoặc làm chính xác hơn: predict cho từng frame).
  \item Encode 30-frame sequence như 1 đoạn video thực.
\end{enumerate}

\textbf{Codec config kèm theo} (bật đồng thời, không bật riêng lẻ):
\begin{itemize}[leftmargin=2em]
  \item \texttt{--framerate 17} — đúng FPS của \texttt{leftImg8bit\_sequence}
        (camera xe ghi $\sim$17 fps).
  \item \texttt{--keyint 0} — để codec tự chọn GOP ($\sim$250), tận dụng
        inter-prediction xuyên suốt 30 frame liên tục.
\end{itemize}

Đây là kịch bản đúng với paper, và là nơi SAC \emph{phát huy} ưu thế: codec coi
ROI là vùng quan trọng xuyên suốt 30 frame, B-frame tham chiếu chéo nhẹ-nén
non-ROI sang heavy-nén non-ROI giúp tiết kiệm bitrate đáng kể.

\subsection{Đề xuất 9 — Loss bổ sung: edge-aware Lovász cho ROI}

ROI có biên dài và bất quy tắc (xe, người, biển báo). Cộng thêm:
\begin{equation}
L_{\text{total}} = L_{\text{CE}} + L_{\text{Dice}} + 0.5\, L_{\text{Lov}}(p_{\text{ROI}}, y_{\text{ROI}})
\end{equation}
với $L_{\text{Lov}}$ là Lovász-hinge loss \emph{chỉ trên lớp ROI} (binary). Lovász
tối ưu trực tiếp IoU $\to$ +$0.5$--$1$ điểm iIoU.

\subsection{Tóm tắt lộ trình triển khai}

\begin{table}[H]
\centering
\caption{Lộ trình theo độ ưu tiên + lợi ích kỳ vọng. ``Standalone'' = bật ngay
trên checkpoint hiện có; ``Combo'' = phải đi kèm điều kiện tiên quyết, không
được bật riêng lẻ.}
\begin{tabular}{clccl}
\toprule
\# & Cải thiện & Loại & Công sức & Kỳ vọng tăng (sau nén) \\
\midrule
1 & \texttt{--roi-source gt} (Đề xuất 1)            & Standalone & $\sim$30 phút        & mIoU $+2$--$3$ pt, iIoU $+1$ pt \\
2 & Khớp bitrate \texttt{--match-bitrate} (Đề xuất 6) & Standalone & 1 giờ chạy           & So sánh fair vs paper \\
3 & Drop metric \texttt{miou\_before/drop} (Đề xuất 7) & Standalone & — (đã tích hợp)      & Diagnostic — không tăng số tuyệt đối \\
4 & Label soft + soft-CE (Đề xuất 2)               & Combo: re-train & 1 giờ + 8 giờ train  & mIoU $+0.3$--$0.5$ pt \\
5 & RGB + cosine LR + 80 ep (Đề xuất 3)            & Combo: re-train & 1 giờ + 8 giờ train  & mIoU $+0.5$--$1$ pt \\
6 & Multi-scale head (Đề xuất 5)                   & Combo: re-train & 2 giờ + 6 giờ train  & mIoU $+0.3$ pt \\
7 & Lovász loss (Đề xuất 9)                        & Combo: re-train & 30 phút + train lại  & iIoU $+0.5$ pt \\
8 & GN thay BN (Đề xuất 4)                         & Combo: re-train & 1 giờ + 8 giờ train  & mIoU $+0.2$ pt nếu batch $\le 4$ \\
9 & leftImg8bit\_sequence + fps17/kt0 (Đề xuất 8)  & Combo: new data & 1 ngày (DL + loader) & SA-PSNR $+1$--$2$ dB \\
\bottomrule
\end{tabular}
\end{table}

\section{Đã triển khai cải tiến trong \texttt{evaluate.py}}

Các Đề xuất \emph{standalone} (\#1, \#6, \#7 — không cần re-train) đã được triển
khai trong \texttt{scripts/new\_branch/evaluate.py} kèm chỉnh sửa hỗ trợ trong
\texttt{compression.py}. Tổng cộng:
\begin{itemize}[leftmargin=2em]
  \item 2 flag \textbf{improvement} mặc định an toàn: \texttt{--roi-source},
        \texttt{--match-bitrate}.
  \item 4 cột \textbf{metric mới} đo drop chuẩn xác hơn:
        \texttt{miou\_before, iiou\_before, miou\_drop, iiou\_drop}.
  \item 3 flag \textbf{advanced opt-in} cho tương lai (không cải thiện một mình,
        cần re-train hoặc đổi dataset): \texttt{--no-gray}, \texttt{--framerate},
        \texttt{--keyint}. Có cảnh báo runtime khi bật sai cảnh.
\end{itemize}

\subsection{Các flag CLI mới}

\begin{table}[H]
\centering
\small
\caption{Flag mới của \texttt{evaluate.py} và Đề xuất tương ứng.}
\begin{tabular}{p{3.8cm} p{1.8cm} p{2.3cm} p{6.4cm}}
\toprule
Flag & Mặc định & Loại & Tác dụng \\
\midrule
\texttt{--roi-source \{pred,gt\}} & \texttt{pred} & Improvement &
  \textbf{Đề xuất \#1}. Khi \texttt{gt}: lấy ROI mask thẳng từ
  \texttt{\_gtFine\_4class.png}, qua khối S2 (macroblock 16$\times$16). Loại lỗi
  segmentation khỏi khối nén — so fair với paper-folder. \\
\addlinespace
\texttt{--match-bitrate} & off & Improvement &
  \textbf{Đề xuất \#6}. SA-* binary-search \texttt{crf\_non}$\in[\text{crf\_roi}{+}1,51]$
  trên \texttt{--match-sample-gops} GOP đầu để bytes SAC khớp baseline H.* $\pm 2\%$. \\
\addlinespace
\texttt{--match-sample-gops INT} & \texttt{1} & Improvement &
  Số GOP làm sample khi binary-search. Tăng để chính xác hơn. \\
\midrule
\texttt{--no-gray} & off & Advanced opt-in &
  Bỏ \texttt{\_rgb\_to\_gray3()} khi tạo tensor input. \emph{\textbf{Chỉ dùng nếu
  checkpoint được train với RGB}} (Đề xuất \#3). Checkpoint mặc định
  \texttt{best\_ccnet\_sac.pth} train grayscale → bật flag này làm mIoU tụt từ
  $\sim$80\% xuống $\sim$50\% do distribution shift (xem \S\ref{sec:pitfall}).
  Có cảnh báo runtime. \\
\addlinespace
\texttt{--framerate INT} & \texttt{30} & Advanced opt-in &
  FPS cho \texttt{ffmpeg}. Mặc định 30 khớp slideshow eval. Đặt 17 chỉ khi đã
  chuyển sang \texttt{leftImg8bit\_sequence} thực (Đề xuất \#8). \\
\addlinespace
\texttt{--keyint INT} & \texttt{30} & Advanced opt-in &
  GOP cho codec. Mặc định 30 — mỗi frame slideshow là I-frame. Đặt 0
  (codec tự chọn $\sim$250) chỉ khi đầu vào là video sequence thực. Sai chỗ:
  bytes/frame tăng $\sim$25\% (xem \S\ref{sec:pitfall}). \\
\bottomrule
\end{tabular}
\end{table}

\subsection{Cột metric mới trong \texttt{per\_frame.csv} / \texttt{summary.csv}}

\begin{itemize}[leftmargin=2em]
  \item \texttt{miou\_before, iiou\_before} — segmentation mIoU/iIoU đo trên frame
        gốc \emph{trước khi nén} (Đề xuất \#7). Đây là ``trần'' của mạng seg.
  \item \texttt{miou\_drop, iiou\_drop} — chênh lệch \texttt{before − after}. Đây
        mới là \emph{chỉ số đo trực tiếp tác hại của nén lên perception}, độc lập
        với chất lượng mạng seg. \texttt{miou\_drop} càng nhỏ (gần 0) càng tốt —
        SAC kỳ vọng \texttt{miou\_drop} nhỏ hơn baseline H.* tại cùng bitrate.
\end{itemize}

\subsection{Thay đổi nội bộ ở \texttt{compression.py}}

\begin{itemize}[leftmargin=2em]
  \item \texttt{\_encode/compress\_traditional/compress\_sac}: thêm tham số
        \texttt{keyint: int | None = 30} (\texttt{None} = bỏ cờ
        \texttt{-x264-params/-x265-params}, để codec dùng GOP mặc định
        $\sim$250). Mặc định 30 — phù hợp pipeline slideshow hiện tại.
  \item \texttt{framerate} mặc định giữ \texttt{30} — khớp với checkpoint cũ và
        evaluate cũ. Đổi sang 17 chỉ khi dùng sequence thực.
\end{itemize}

(Lưu ý lịch sử: phiên bản ngắn ngày 18/05 từng đặt defaults \texttt{17/None},
nhưng làm bytes/frame tăng và miou\_before vô lý với slideshow eval — đã đảo lại
sau khi gặp pitfall \S\ref{sec:pitfall}.)

\subsection{Tổ chức lại pipeline trong \texttt{evaluate.py}}

\begin{enumerate}[leftmargin=2em]
  \item \textbf{Tiền xử lý gộp một lượt} (mục [1/3] trong log):
        load mỗi frame, tính \emph{seg-before} (mIoU/iIoU trên frame gốc), từ đó
        \emph{vừa} có baseline mIoU/iIoU, \emph{vừa} có ROI mask khi
        \texttt{--roi-source pred}. Không cần predict 2 lần.
  \item \textbf{Sắp lại thứ tự method} (mục [2/3]): chạy H.264/H.265 baseline
        trước (sorted by \texttt{startswith("SA-")}), lưu \texttt{bytes\_per\_frame}
        vào dict \texttt{baseline\_bytes}. Khi tới SA-* mà có
        \texttt{--match-bitrate}, lấy target ra để binary-search.
  \item \textbf{Helper \texttt{\_frame\_to\_seg\_tensor}}: gom logic
        ``resize → grayscale (optional) → normalize → tensor''. Tránh import
        \texttt{\_normalize/\_rgb\_to\_gray3} bên trong hot-loop.
  \item \textbf{Lưu \texttt{config.json}}: ghi lại toàn bộ tham số eval (split,
        flags, ckpt) cùng output để reproducible khi đối chiếu giữa nhiều lần chạy.
\end{enumerate}

\subsection{Pitfall đã gặp khi test 2026-05-18 — và cách khắc phục}
\label{sec:pitfall}

Khi chạy đầu tiên với defaults \texttt{--framerate 17 --keyint 0 --no-gray}, kết
quả thua hẳn run cũ. Sau khi đảo defaults về \texttt{30/30/gray-on}, kết quả phục
hồi và Đề xuất \#1 thực sự cho cải thiện.

\begin{table}[H]
\centering
\small
\caption{Chuỗi 3 run: run cũ (chưa có cải tiến) $\to$ run lỗi (bật sai Đề xuất
\#2/\#7) $\to$ run phục hồi (defaults an toàn). Đơn vị mIoU/iIoU = \%.}
\begin{tabular}{lcccc}
\toprule
Run & H264 mIoU after & SA-X264 mIoU after & H264 bytes/frame & mIoU \emph{before} \\
\midrule
20260517\_204131 (100 fr, pred, gray, kt=30, fps=30) & 76.42 & 75.31 & 59.5K  & — \\
20260518\_010212 (10 fr, gt, RGB, kt=0, fps=17)      & 50.65 & 48.74 & 73.2K  & \textbf{48.4} \\
20260518\_011924 (10 fr, pred, gray, kt=30, fps=30)  & 80.21 & 80.01 & 64.1K  & \textbf{82.23} \\
\bottomrule
\end{tabular}
\end{table}

\textbf{Hai nguyên nhân khiến run thứ 2 thất bại:}

\paragraph{(a) \texttt{--no-gray} mismatch với checkpoint.} Dataset training
(\texttt{dataset.py:100}) \emph{luôn} convert grayscale trước khi normalize. Mọi
checkpoint hiện có đều thấy input grayscale lúc training. Khi inference với
RGB, model nhận input out-of-distribution → predict gần như random ở các class
phụ thuộc chroma. Bằng chứng: \texttt{miou\_before} chỉ 48.4\% trên frame gốc
(chưa nén), và \texttt{miou\_drop} \emph{âm tất cả} (seg trên frame nén tốt hơn
frame gốc — vô lý, dấu hiệu model đang trả lời nhiễu).

\paragraph{(b) \texttt{keyint=0} sai với slideshow eval.} Test split của
new-branch hiện không phải video thực — mỗi frame thuộc 1 thành phố/scene khác
nhau. Khi đặt \texttt{keyint=0} (codec dùng GOP $\sim$250), x264 cố làm P-frame
qua scene cut → fail, fall back sang noisy P-frames hoặc force I-frame ngầm →
bytes/frame H264 tăng từ 59KB → 73KB ($+23\%$).

\paragraph{Khắc phục (run thứ 3, 20260518\_011924).} Đảo defaults về
\texttt{--no-gray off} (giữ grayscale) + \texttt{--keyint 30} +
\texttt{--framerate 30} + sanity check runtime cảnh báo khi
\texttt{miou\_before $< 60\%$}. Kết quả phục hồi:
\begin{itemize}[leftmargin=2em]
  \item \texttt{miou\_before} = 82.23\% (so với 48.4\% lần lỗi) — model trả lời
        đúng cảnh.
  \item \texttt{miou\_drop} $= +2.0\%$ (H264), $+2.2\%$ (SA-X264) — dương
        (đúng chiều: nén làm seg tệ đi một chút).
  \item bytes/frame H264 = 64.1K (so với 73.2K lần lỗi) — codec không lãng phí
        I-frame ngầm.
  \item SA-X264 SA-PSNR (39.30) $>$ H264 (39.23) — SAC win nhẹ kể cả với
        bytes/frame cao hơn (90.7K vs 64.1K, do chưa bật
        \texttt{--match-bitrate}).
\end{itemize}

\paragraph{Bài học.} Các Đề xuất \emph{thay đổi input pipeline hoặc codec
behavior} là \textbf{combo có điều kiện tiên quyết}:
\begin{itemize}[leftmargin=2em]
  \item Bỏ grayscale (Đề xuất \#3) ⇒ phải re-train từ đầu.
  \item FPS=17 + \texttt{keyint=0} (Đề xuất \#8) ⇒ phải có
        \texttt{leftImg8bit\_sequence} thực.
\end{itemize}
Bật riêng lẻ không re-train hoặc không đổi dataset là \emph{phản tác dụng}. Vì
vậy default của \texttt{evaluate.py} giữ giá trị \emph{tương thích checkpoint
hiện có}; flag advanced phải bật thủ công sau khi đã thoả prerequisite.

\subsection{Cách dùng — các kịch bản khuyến nghị}

\paragraph{Kịch bản A — so trực tiếp với paper-folder (oracle mask):}
\begin{lstlisting}[language=bash]
# Dung GT mask, GIU grayscale (khop checkpoint hien tai).
python -m scripts.new_branch.evaluate \
    --split test --roi-source gt
\end{lstlisting}
Kỳ vọng: mIoU sau nén $\approx 76$--$80$\%, iIoU $\approx 91$--$93$\% với checkpoint
\texttt{best\_ccnet\_sac.pth} hiện có (đã verify run 20260518\_011924: \texttt{miou\_before}
= 82.23\%, mIoU after = 80.21\% / 80.01\% / 79.48\% / 78.40\% cho H264/SA-X264/H265/SA-X265).
Để đạt $\approx 90\%$/$92\%$ như paper Bảng IV cần checkpoint mạnh hơn — xem Đề
xuất \#2 (label soft), \#3 (RGB+cosine), \#4 (GN), \#5 (multi-scale) ở mục
``Đề xuất chưa triển khai''.

\paragraph{Kịch bản B — end-to-end honest (pred mask) + khớp bitrate:}
\begin{lstlisting}[language=bash]
python -m scripts.new_branch.evaluate \
    --split test --roi-source pred \
    --match-bitrate --match-sample-gops 2
\end{lstlisting}
Cùng bitrate với baseline H.*; đo \texttt{miou\_drop}/\texttt{iiou\_drop} giữa SAC
và H.* — nếu SAC drop nhỏ hơn thì SAC thắng.

\paragraph{Kịch bản C — debug nhanh trên 20 frame:}
\begin{lstlisting}[language=bash]
python -m scripts.new_branch.evaluate \
    --split test --num-frames 20 --gop 10 \
    --roi-source gt
\end{lstlisting}

\paragraph{Kịch bản D — sau khi re-train với RGB + leftImg8bit\_sequence (advanced):}
\begin{lstlisting}[language=bash]
python -m scripts.new_branch.evaluate \
    --split test --roi-source gt \
    --no-gray --framerate 17 --keyint 0
\end{lstlisting}
Chỉ chạy khi đã: (1) sửa \texttt{dataset.py} bỏ \texttt{\_rgb\_to\_gray3}, re-train
$\to$ checkpoint mới; (2) thay loader sang sequence thực 30 frame liên tục.

\subsection{Đề xuất chưa triển khai (cần re-train hoặc dữ liệu mới)}

\begin{itemize}[leftmargin=2em]
  \item \textbf{Đề xuất \#2 (label resize mềm)} — sửa \texttt{dataset.py} và
        \texttt{losses.py}, ảnh hưởng train.
  \item \textbf{Đề xuất \#3 (RGB + cosine LR + augment mạnh + 80 epoch)} —
        \texttt{dataset.py} (bỏ grayscale) + \texttt{train.py}.
  \item \textbf{Đề xuất \#4 (GN thay BN)} — sửa
        \texttt{scripts/CCNet/networks/ccnet.py} + re-train.
  \item \textbf{Đề xuất \#5 (multi-scale head)} — sửa kiến trúc + re-train.
  \item \textbf{Đề xuất \#8 (\texttt{leftImg8bit\_sequence} thực + fps17/kt0)} —
        cần download 324GB + viết loader mới (\texttt{dataset.py} hiện chỉ load
        frame có nhãn).
  \item \textbf{Đề xuất \#9 (Lovász loss)} — chỉ \texttt{losses.py} + re-train.
\end{itemize}

\section{Kết luận}

paper-folder \textbf{không phải} là tái hiện 100\% trung thành với mô tả của Wang et
al. 2023 — kiến trúc backbone của nó là CNN tự thiết kế nhỏ với GroupNorm, không phải
ResNet-101 dilated như paper viết. Tuy nhiên paper-folder lại đạt kết quả tốt vì:
\begin{enumerate}[leftmargin=2em]
  \item \textbf{Dùng GT mask để tạo two-stream — đây là điểm cốt yếu.} Bảng IV của
        paper khi báo cáo iIoU $=92.45\%$ ngầm giả định pipeline có khối S1 đủ tốt;
        paper-folder ``cheat'' bằng cách dùng thẳng GT, bỏ qua sai số khối S1.
  \item Multi-scale prediction fusion + GroupNorm + label resize mềm — ba kỹ thuật
        nhỏ này cộng lại tạo lợi thế đáng kể trên 4 lớp gộp.
  \item Train from scratch tránh được ``BN-shift'' trong vài epoch đầu mà new-branch
        gặp khi fine-tune ResNet-101 ImageNet.
\end{enumerate}

new-branch theo đúng spec paper (ResNet-101 + RCCA, R=2, DSN) là \textbf{phương pháp
luận đúng}, nhưng phải bổ sung 3 việc tối thiểu để bằng kết quả:
\begin{enumerate}[leftmargin=2em]
  \item Thêm chế độ \texttt{--roi-source gt} để so trực tiếp với paper-folder.
  \item Bỏ grayscale, giữ RGB.
  \item Sửa codec config về FPS=17 + bỏ ép \texttt{keyint=30}.
\end{enumerate}

\textbf{Đề xuất \#1 đã được triển khai an toàn} trong \texttt{evaluate.py}
(\texttt{--roi-source gt}), cộng với cờ \texttt{--match-bitrate} (Đề xuất \#6) và
các cột metric \texttt{miou\_before/iiou\_before/miou\_drop/iiou\_drop}
(Đề xuất \#7). New-branch giờ có thể \emph{so fair} với paper-folder \emph{ở hai chế
độ}:
\begin{itemize}[leftmargin=2em]
  \item \emph{Oracle} (\texttt{--roi-source gt}) — kỳ vọng bằng paper-folder.
  \item \emph{End-to-end} (\texttt{--roi-source pred} + \texttt{--match-bitrate})
        — đánh giá đúng tác dụng SAC trong deployment.
\end{itemize}

Riêng việc ``bỏ grayscale'' và ``đổi codec config (fps=17, keyint=0)'' \textbf{đã
được rút khỏi danh sách improvement standalone}: chúng là combo có điều kiện tiên
quyết (cần re-train hoặc cần \texttt{leftImg8bit\_sequence}). Bằng chứng cụ thể:
khi bật riêng lẻ, mIoU sụt từ 76\% xuống 50\% và bytes/frame tăng 23\%
(\S\ref{sec:pitfall}). Vì vậy giờ chúng được giữ làm \emph{advanced opt-in
flags} kèm cảnh báo runtime, và phần ``re-train RGB'' đã được gộp vào Đề xuất
\#3, phần ``fps=17 + keyint=0'' gộp vào Đề xuất \#8.

Các Đề xuất 2--5, 8, 9 còn lại đòi hỏi re-train hoặc dữ liệu mới — dành cho ai
muốn \emph{vượt} paper-folder (mIoU $\ge 91\%$, iIoU $\ge 93\%$) trong khuôn khổ
ResNet-101 chuẩn.

\bigskip
\noindent\textbf{File so sánh trực quan:}
\begin{itemize}[leftmargin=2em]
  \item paper-folder: \texttt{scripts/Semantic-Aware-Video-Compression-for-Automotive-Cameras/Codes/CCNet/\{CCNet.py, train.py, TwoStream\_generate.py\}}
  \item new-branch: \texttt{scripts/new\_branch/\{model.py, train.py, evaluate.py, stream\_separation.py\}}
\end{itemize}

\end{document}
